\chapter{Stolz–Cesàro theorem}

\begin{lemma}[Addition of Strict Inequalities]
    \label{lem:Addition_of_Strict_Inequalities_Property}
    Let $A, B, C, D$ be real numbers. If $A < B$ and $C < D$, then $A + C < B + D$.
\end{lemma}

\begin{proof}
    The inequality $A < B$ is equivalent to $B - A > 0$. Similarly, $C < D$ is equivalent to $D - C > 0$.
    The sum of two positive real numbers is positive, so $(B - A) + (D - C) > 0$.
    By rearranging the terms, we get $(B + D) - (A + C) > 0$.
    This is equivalent to the inequality $A + C < B + D$.
\end{proof}

\begin{lemma}[Multiplication of an Inequality by a Positive Number]
    \label{lem:Multiplication_by_Positive}
    Let $x, y, z$ be real numbers. If $x < y$ and $z > 0$, then $xz < yz$.
\end{lemma}

\begin{proof}
    The inequality $x < y$ is equivalent to $y - x > 0$.
    We are given that $z > 0$. The product of two positive real numbers is positive, so $(y - x)z > 0$.
    By the distributive property, this is $yz - xz > 0$.
    This is equivalent to the inequality $xz < yz$.
\end{proof}

\begin{lemma}[Multiplication of an Inequality by a Negative Number]
    \label{lem:Multiplication_by_Negative}
    Let $x, y, z$ be real numbers. If $x < y$ and $z < 0$, then $xz > yz$.
\end{lemma}

\begin{proof}
    The inequality $x < y$ is equivalent to $y - x > 0$.
    We are given that $z < 0$. The product of a positive real number and a negative real number is negative, so $(y - x)z < 0$.
    By the distributive property, this is $yz - xz < 0$.
    This is equivalent to the inequality $yz < xz$, or $xz > yz$.
\end{proof}

\begin{lemma}[Additivity of Strict Inequalities over a Sum]
    \label{lem:Additivity_of_Strict_Inequalities}
    \uses{lem:Addition_of_Strict_Inequalities_Property}
    Let $(x_k)$ and $(y_k)$ be two sequences of real numbers. If $x_k < y_k$ for all integers $k$ such that $M \le k \le N$, then $\sum_{k=M}^{N} x_k < \sum_{k=M}^{N} y_k$.
\end{lemma}

\begin{proof}
    We prove this by induction on the number of terms, $n = N - M + 1$.
    Base case: For $n=1$, we have $N=M$. The statement is $x_M < y_M$, which is true by hypothesis.
    Inductive step: Assume the statement is true for $n$ terms. Consider a sum over $n+1$ terms, from $M$ to $M+n$. We have:
    $$\sum_{k=M}^{M+n} x_k = \left(\sum_{k=M}^{M+n-1} x_k\right) + x_{M+n}$$
    By the inductive hypothesis, $\sum_{k=M}^{M+n-1} x_k < \sum_{k=M}^{M+n-1} y_k$. By the given condition, $x_{M+n} < y_{M+n}$.
    Using lemma \ref{lem:Addition_of_Strict_Inequalities_Property}, we can add these two inequalities:
    $$\left(\sum_{k=M}^{M+n-1} x_k\right) + x_{M+n} < \left(\sum_{k=M}^{M+n-1} y_k\right) + y_{M+n}$$
    This simplifies to $\sum_{k=M}^{M+n} x_k < \sum_{k=M}^{M+n} y_k$.
    The statement holds for $n+1$ terms. By the principle of mathematical induction, the lemma is true.
\end{proof}

\begin{lemma}[Telescoping Sum]
    \label{lem:Telescoping_Sum}
    For any sequence of real numbers $(x_k)$ and integers $n > M$, the following identity holds: $\sum_{k=M}^{n-1} (x_{k+1} - x_k) = x_n - x_M$.
\end{lemma}

\begin{proof}
    We prove this by induction on $n$.
    Base case: Let $n=M+1$. The sum is $\sum_{k=M}^{M} (x_{k+1} - x_k) = x_{M+1} - x_M$. The right hand side is $x_{M+1} - x_M$. The identity holds.
    Inductive step: Assume the identity holds for some $n > M$. We want to show it holds for $n+1$.
    $$\sum_{k=M}^{n} (x_{k+1} - x_k) = \left(\sum_{k=M}^{n-1} (x_{k+1} - x_k)\right) + (x_{n+1} - x_n)$$
    By the inductive hypothesis, the sum on the right is equal to $x_n - x_M$.
    $$= (x_n - x_M) + (x_{n+1} - x_n) = x_{n+1} - x_M$$
    This is the required formula for $n+1$. By the principle of mathematical induction, the lemma is proven.
\end{proof}

\begin{lemma}[Property of Sequences Diverging to Infinity]
    \label{lem:Property_of_Divergent_Sequences}
    If a sequence $(b_n)$ diverges to $+\infty$, then for any real number $M$, there exists an integer $N$ such that $b_n > M$ for all $n > N$. In particular, the sequence is eventually positive.
\end{lemma}

\begin{proof}
    The statement is the definition of divergence to $+\infty$. For any $M \in \mathbb{R}$, the definition of $\lim_{n \to \infty} b_n = +\infty$ states that there exists an integer $N$ such that for all $n > N$, we have $b_n > M$. To prove that the sequence is eventually positive, we can choose $M=0$. Then there exists an $N$ such that for all $n > N$, $b_n > 0$.
\end{proof}

\begin{lemma}[Property of Monotone Sequences Converging to Zero]
    \label{lem:Property_of_Monotone_Sequences_Converging_to_Zero}
    If a sequence $(b_n)$ is strictly decreasing and $\lim_{n\to\infty} b_n = 0$, then $b_n > 0$ for all $n$.
\end{lemma}

\begin{proof}
    Assume for contradiction that there exists an integer $N$ such that $b_N \leq 0$. Since the sequence $(b_n)$ is strictly decreasing, for all $n > N$, we have $b_n < b_N \leq 0$. This implies that the limit of the sequence, if it exists, must be less than or equal to $b_N$, so $\lim_{n\to\infty} b_n \leq b_N \leq 0$. However, we are given that the limit is 0. This requires $b_N = 0$. But if $b_N = 0$, then $b_{N+1} < b_N = 0$, which means $\lim_{n\to\infty} b_n \leq b_{N+1} < 0$. This contradicts the fact that the limit is 0. Therefore, our initial assumption must be false, and $b_n > 0$ for all $n$.
\end{proof}

\begin{lemma}[Preservation of Non-Strict Inequalities under Limits]
    \label{lem:Limit_Inequality_Property}
    Let $(x_n)$ and $(y_n)$ be convergent sequences. If there exists an integer $N$ such that $x_n \leq y_n$ for all $n > N$, then $\lim_{n \to \infty} x_n \leq \lim_{n \to \infty} y_n$.
\end{lemma}

\begin{proof}
    Let $L_x = \lim_{n \to \infty} x_n$ and $L_y = \lim_{n \to \infty} y_n$. Assume for contradiction that $L_x > L_y$. Let $\epsilon = \frac{L_x - L_y}{2}$. Since $\epsilon > 0$, by the definition of a limit, there exist integers $N_x$ and $N_y$ such that:
    For $n > N_x$, $|x_n - L_x| < \epsilon \implies x_n > L_x - \epsilon$.
    For $n > N_y$, $|y_n - L_y| < \epsilon \implies y_n < L_y + \epsilon$.
    For any $n > \max(N, N_x, N_y)$, we have both conditions satisfied. Substituting the definition of $\epsilon$:
    $x_n > L_x - \frac{L_x - L_y}{2} = \frac{L_x + L_y}{2}$.
    $y_n < L_y + \frac{L_x - L_y}{2} = \frac{L_x + L_y}{2}$.
    This implies $y_n < \frac{L_x + L_y}{2} < x_n$, so $y_n < x_n$. This contradicts the hypothesis that $x_n \leq y_n$ for all $n>N$. Thus, our assumption must be false, and we must have $L_x \leq L_y$.
\end{proof>

\begin{lemma}[Bounding the sequence difference]
    \label{lem:Bounding_the_sequence_difference}
    \uses{lem:Multiplication_by_Positive}
    \uses{lem:Additivity_of_Strict_Inequalities}
    \uses{lem:Telescoping_Sum}
    Let $(a_n)_{n\geq 1}$ and $(b_n)_{n\geq 1}$ be two sequences of real numbers. Assume that $(b_n)_{n\geq 1}$ is a strictly increasing sequence such that $b_n \to +\infty$.
    If $\lim_{n\to\infty} \frac{a_{n+1}-a_n}{b_{n+1}-b_n} = l$, where $l$ is a finite real number, then for any $\epsilon > 0$, there exists an integer $\nu$ such that for all $n > \nu+1$, the following inequality holds:
    $$(l - \epsilon/2)(b_n - b_{\nu+1}) + a_{\nu+1} < a_n < (l + \epsilon/2)(b_n - b_{\nu+1}) + a_{\nu+1}$$
\end{lemma}

\begin{proof}
    Given $\lim_{n\to\infty} \frac{a_{n+1}-a_n}{b_{n+1}-b_n} = l$, for any $\epsilon/2 > 0$, there exists an integer $\nu$ such that for all $k > \nu$:
    $$l - \frac{\epsilon}{2} < \frac{a_{k+1}-a_k}{b_{k+1}-b_k} < l + \frac{\epsilon}{2}$$
    Since $(b_n)$ is strictly increasing, $b_{k+1} - b_k > 0$. Using lemma \ref{lem:Multiplication_by_Positive}, we can multiply by the positive term $b_{k+1}-b_k$ without changing the direction of the inequalities:
    $$(l - \frac{\epsilon}{2})(b_{k+1}-b_k) < a_{k+1}-a_k < (l + \frac{\epsilon}{2})(b_{k+1}-b_k)$$
    This holds for all integers $k$ from $\nu+1$ to $n-1$. Applying lemma \ref{lem:Additivity_of_Strict_Inequalities}, we sum the inequalities:
    $$\sum_{k=\nu+1}^{n-1} (l - \frac{\epsilon}{2})(b_{k+1}-b_k) < \sum_{k=\nu+1}^{n-1} (a_{k+1}-a_k) < \sum_{k=\nu+1}^{n-1} (l + \frac{\epsilon}{2})(b_{k+1}-b_k)$$
    Using lemma \ref{lem:Telescoping_Sum} on the sums involving $a_k$ and $b_k$, we get:
    $$(l - \frac{\epsilon}{2})(b_n - b_{\nu+1}) < a_n - a_{\nu+1} < (l + \frac{\epsilon}{2})(b_n - b_{\nu+1})$$
    Adding $a_{\nu+1}$ to all parts of the inequality gives the desired result.
\end{proof}

\begin{lemma}[Vanishing of the residual term]
    \label{lem:Vanishing_of_the_residual_term}
    \uses{lem:Property_of_Divergent_Sequences}
    Let $(b_n)_{n\geq 1}$ be a sequence such that $b_n \to +\infty$. For any constant $C$, we have:
    $$\lim_{n\to\infty} \frac{C}{b_n} = 0$$
\end{lemma}
\begin{proof}
    Given any $\epsilon > 0$, we need to find an integer $N$ such that for all $n > N$, $\left|\frac{C}{b_n} - 0\right| < \epsilon$.
    This inequality is equivalent to $\frac{|C|}{|b_n|} < \epsilon$, which simplifies to $|b_n| > \frac{|C|}{\epsilon}$.
    Since $b_n \to +\infty$, by lemma \ref{lem:Property_of_Divergent_Sequences}, for any real number $M$, there exists an $N_1$ such that $b_n > M$ for all $n > N_1$. Let's choose $M = \frac{|C|}{\epsilon}$.
    Also by lemma \ref{lem:Property_of_Divergent_Sequences}, there exists $N_2$ such that $b_n > 0$ for $n > N_2$.
    Let $N = \max(N_1, N_2)$. For $n > N$, we have $b_n > \frac{|C|}{\epsilon}$ and $b_n>0$. Thus $|b_n| = b_n > \frac{|C|}{\epsilon}$.
    This satisfies the condition for the limit to be 0.
\end{proof}

\begin{theorem}[Stolz–Cesàro Theorem for the */∞ case]
    \label{thm:Stolz-Cesaro_Theorem_for_the_infinity_case}
    \uses{lem:Bounding_the_sequence_difference}
    \uses{lem:Vanishing_of_the_residual_term}
    \uses{lem:Property_of_Divergent_Sequences}
    Let $(a_n)_{n\geq 1}$ and $(b_n)_{n\geq 1}$ be two sequences of real numbers. Assume that $(b_n)_{n\geq 1}$ is a strictly monotone and divergent sequence. If the following limit exists:
    $$\lim_{n\to\infty} \frac{a_{n+1}-a_n}{b_{n+1}-b_n} = l$$
    Then, the limit $\lim_{n\to\infty} \frac{a_n}{b_n}$ also exists and is equal to $l$.
\end{theorem}

\begin{proof}
    We prove the case where $(b_n)$ is strictly increasing and $b_n \to +\infty$. The case where it is strictly decreasing and tends to $-\infty$ is similar.
    From lemma \ref{lem:Bounding_the_sequence_difference}, for any $\epsilon > 0$, there exists $\nu$ such that for $n > \nu+1$:
    $$(l - \epsilon/2)(b_n - b_{\nu+1}) + a_{\nu+1} < a_n < (l + \epsilon/2)(b_n - b_{\nu+1}) + a_{\nu+1}$$
    Since $b_n \to +\infty$, by lemma \ref{lem:Property_of_Divergent_Sequences}, there exists an $n_0$ such that $b_n > 0$ for all $n > n_0$. For $n > \max(\nu+1, n_0)$, we can divide the inequality by $b_n$:
    $$(l - \epsilon/2)\left(1 - \frac{b_{\nu+1}}{b_n}\right) + \frac{a_{\nu+1}}{b_n} < \frac{a_n}{b_n} < (l + \epsilon/2)\left(1 - \frac{b_{\nu+1}}{b_n}\right) + \frac{a_{\nu+1}}{b_n}$$
    This can be rewritten as:
    $$(l - \epsilon/2) + \frac{a_{\nu+1} - b_{\nu+1}(l - \epsilon/2)}{b_n} < \frac{a_n}{b_n} < (l + \epsilon/2) + \frac{a_{\nu+1} - b_{\nu+1}(l + \epsilon/2)}{b_n}$$
    The terms $\frac{a_{\nu+1} - b_{\nu+1}(l \pm \epsilon/2)}{b_n}$ are of the form $C/b_n$. By lemma \ref{lem:Vanishing_of_the_residual_term}, these terms go to 0 as $n \to \infty$.
    Thus, there exist $N_1$ and $N_2$ such that for $n>N_1$, $\left|\frac{a_{\nu+1} - b_{\nu+1}(l - \epsilon/2)}{b_n}\right| < \epsilon/2$, and for $n>N_2$, $\left|\frac{a_{\nu+1} - b_{\nu+1}(l + \epsilon/2)}{b_n}\right| < \epsilon/2$.
    For $n > \max(\nu+1, n_0, N_1, N_2)$, we have:
    $$l - \epsilon = (l - \epsilon/2) - \epsilon/2 < \frac{a_n}{b_n} < (l + \epsilon/2) + \epsilon/2 = l + \epsilon$$
    This shows that for any $\epsilon > 0$, there is an $N$ such that for $n>N$, $|\frac{a_n}{b_n} - l| < \epsilon$. Therefore, $\lim_{n\to\infty} \frac{a_n}{b_n} = l$.
\end{proof}

\begin{theorem}[Stolz–Cesàro Theorem for the 0/0 case]
    \label{thm:Stolz-Cesaro_Theorem_for_the_0_case}
    \uses{lem:Multiplication_by_Negative}
    \uses{lem:Additivity_of_Strict_Inequalities}
    \uses{lem:Telescoping_Sum}
    \uses{lem:Limit_Inequality_Property}
    \uses{lem:Property_of_Monotone_Sequences_Converging_to_Zero}
    Let $(a_n)_{n\geq 1}$ and $(b_n)_{n\geq 1}$ be two sequences of real numbers such that $\lim_{n\to\infty} a_n = 0$ and $\lim_{n\to\infty} b_n = 0$. Assume $(b_n)_{n\geq 1}$ is strictly decreasing. If
    $$\lim_{n\to\infty} \frac{a_{n+1}-a_n}{b_{n+1}-b_n} = l$$
    then
    $$\lim_{n\to\infty} \frac{a_n}{b_n} = l$$
\end{theorem}
\begin{proof}
    For any $\epsilon > 0$, there is an integer $N$ such that for all $k \geq N$,
    $$l - \epsilon < \frac{a_{k+1}-a_k}{b_{k+1}-b_k} < l + \epsilon$$
    Since $(b_n)$ is strictly decreasing, $b_{k+1}-b_k < 0$. Using lemma \ref{lem:Multiplication_by_Negative}, we multiply by the negative term $b_{k+1}-b_k$, which reverses the inequalities:
    $$(l+\epsilon)(b_{k+1}-b_k) < a_{k+1}-a_k < (l-\epsilon)(b_{k+1}-b_k)$$
    Let $n > N$. For any $m > n$, we sum from $k=n$ to $m-1$. By lemma \ref{lem:Additivity_of_Strict_Inequalities}:
    $$\sum_{k=n}^{m-1} (l+\epsilon)(b_{k+1}-b_k) < \sum_{k=n}^{m-1} (a_{k+1}-a_k) < \sum_{k=n}^{m-1} (l-\epsilon)(b_{k+1}-b_k)$$
    Using lemma \ref{lem:Telescoping_Sum}, we simplify the sums:
    $$(l+\epsilon)(b_m-b_n) < a_m-a_n < (l-\epsilon)(b_m-b_n)$$
    Now we take the limit as $m \to \infty$. Since $\lim_{m\to\infty} a_m = 0$ and $\lim_{m\to\infty} b_m = 0$, by lemma \ref{lem:Limit_Inequality_Property}, we can take the limit inside the inequalities:
    $$(l+\epsilon)(-b_n) \leq -a_n \leq (l-\epsilon)(-b_n)$$
    By lemma \ref{lem:Property_of_Monotone_Sequences_Converging_to_Zero}, since $(b_n)$ is strictly decreasing to 0, we have $b_n > 0$ for all $n$. Thus, $-b_n < 0$. Dividing the entire inequality by $-b_n$ reverses the inequalities again:
    $$l-\epsilon \leq \frac{a_n}{b_n} \leq l+\epsilon$$
    This can be written as $|\frac{a_n}{b_n} - l| \leq \epsilon$. Since this holds for any $n>N$ and for any $\epsilon > 0$, we conclude that $\lim_{n\to\infty} \frac{a_n}{b_n} = l$.
\end{proof}